% Work-MAT | 2. rocnik | pripravovana hodina 45
% Návaznost: uživatelské připravy.zip (hodiny 01–20), autorská učebnice na GitHubu.
% Realisticky.cz: didaktická inspirace v posloupnosti a typech úloh; zadání jsou vlastní.
\documentclass[11pt,a4paper]{article}
\input{../../sablony/pripravy/priprava.sty}
\nastavPripravu{2. ročník}{Opakování funkcí a rovnic}{45}{Opakování exponenciál a logaritmů}
\begin{document}
\blok{Motivace a návaznost}
Propojuj obě navzájem inverzní operace. První úloha je exponenciální, druhá logaritmická; v závěru žáci sami formulují, kdy přechod využijí.
\poznamka{Orientační nabídka pro 45--55 minut: krátké zopakování a motivace, nový výklad, tři společné vzorové úlohy, samostatná práce. Vyber podle tempa třídy; zbývající zadání slouží jako rezerva.}
\blok{Výklad}
Pro $a>0$, $a\ne1$ a $b>0$ platí $a^{\log_a b}=b$ a $\log_a(a^x)=x$. V podmínkách logaritmických úloh nesmí vzniknout nepřípustný argument.
\blok{Společné příklady}
\begin{prikladReseniBox}[title={Příklad 1}]
Vyřeš, resp. urči: $3^x=27$.
\tcblower
\[\begin{aligned}3^x=3^3 \\ x=3\end{aligned}\]
\end{prikladReseniBox}
\begin{prikladReseniBox}[title={Příklad 2}]
Vyřeš, resp. urči: $\log_3(x-1)=2$.
\tcblower
\[\begin{aligned}x>1 \\ x-1=9 \\ x=10\end{aligned}\]
\end{prikladReseniBox}
\poznamka{Nezaváděj metody výpočtu obecných logaritmů pomocí kalkulátoru jako náhradu přesných známých mocnin.}
\blok{Samostatné procvičování}
\begin{zadaniReseniBox}[title={Samostatná práce 1}]
\begin{enumerate}[label=\alph*)]
\item $2^x=2$.
\item $2^x=4$.
\item $2^x=8$.
\item $2^x=16$.
\item $2^x=32$.
\item $3^x=3$.
\end{enumerate}
\tcblower
\begin{enumerate}[label=\alph*)]
\item $x=1$.
\item $x=2$.
\item $x=3$.
\item $x=4$.
\item $x=5$.
\item $x=1$.
\end{enumerate}
\end{zadaniReseniBox}
\blok{Dobrovolné rozšíření}
\begin{zadaniReseniBox}[title={Rozšíření 1}]
\begin{enumerate}[label=\alph*)]
\item $Vytvoř vlastní úlohu tohoto typu a se spolužákem porovnej podmínky a výsledek.$.
\end{enumerate}
\tcblower
\begin{enumerate}[label=\alph*)]
\item Zkontroluj definice a argumenty..
\end{enumerate}
\end{zadaniReseniBox}
\blok{Úkol na doma}
U každého logaritmu vypiš podmínky argumentu a proveď zkoušku. U rovnic vyznač případně nepřípustné kořeny, u nerovnic odpověz intervalem.
\begin{zadaniReseniBox}[title={Domácí úkol}]
\begin{enumerate}[label=\alph*)]
\item $3^x=9$.
\item $3^x=27$.
\item $3^x=81$.
\item $3^x=243$.
\item $5^x=5$.
\item $5^x=25$.
\item $5^x=125$.
\item $5^x=625$.
\end{enumerate}
\tcblower
\begin{enumerate}[label=\alph*)]
\item $x=2$.
\item $x=3$.
\item $x=4$.
\item $x=5$.
\item $x=1$.
\item $x=2$.
\item $x=3$.
\item $x=4$.
\end{enumerate}
\end{zadaniReseniBox}
\poznamka{Domácí úkol obsahuje 8 položek; předpokládaná domácí práce přibližně 15--30 minut. Po větších blocích jsou vědomě promíchané dříve probrané typy.}
\end{document}
